\chapter{Higher-order programming}\label{ch:hoprog}

% local notation: the sequent of Figure 6.5 that processes a multiset of E-formulas
\newcommand{\hoprogE}[4]{#1;#2\mathrel{-\!\langle #3\rangle\!\!\longrightarrow}#4}
\newcommand{\hoprogT}[1]{\lBrack #1\rBrack}

\Cref{ch:lambda} replaced first-order terms by simply typed $\lambda$-terms and allowed quantification at
every type, including predicate types (types whose target is $\typ{o}$); it also showed that unification of
such terms is undecidable and may have infinitely many incomparable solutions. This chapter, which follows
Chapter~5 of the book \cite{MillerNadathur2012} and the core of its Chapter~6, asks the practical
question: \emph{what logic programming languages do we get, and what do they buy us?}

Two kinds of higher-order variable will appear, and it is worth keeping them apart from the start.
\begin{itemize}
\item \textbf{Predicate variables} (target type $\typ{o}$) are instantiated by relations and then
  \emph{invoked as goals}. They give higher-order programming in the style a Scheme programmer knows:
  \lp{mappred}, folds, closures of relations, continuation passing.
\item \textbf{Non-predicate function variables} (target type not $\typ{o}$) are instantiated by
  $\lambda$-terms that are only ever \emph{evaluated by $\beta\eta$-conversion}. Their values can also be
  \emph{found by solving unification problems}; the book calls this ``perhaps the single most novel aspect
  of higher-order logic programming'' \book{138}.
\end{itemize}
We first fix which formulas may be goals and which may be clauses (\cref{sec:hoprog-languages}), then
program with predicate variables (\cref{sec:hoprog-hoprog,sec:hoprog-flex,sec:hoprog-reason,sec:hoprog-logconst,sec:hoprog-naf}),
then with function variables (\cref{sec:hoprog-functions,sec:hoprog-panacea}), compare with functional
programming (\cref{sec:hoprog-fp}), and finally show how the same two logical devices that hide a
predicate, $\forall$ in goals and $\imp$ in goals, also give a modules language
(\cref{sec:hoprog-modules}).

A question runs through the whole chapter, and the kanren boxes answer it case by case: \emph{in a
miniKanren hosted in Scheme, where a relation is a Scheme procedure, which of these features are
free, and which really need $\lambda$-terms in the term language?} \Cref{sec:hoprog-free} collects the
answers.

\begin{tutorial}[How to run the Scheme listings]
Every Scheme listing in this chapter was run on today's \micro{} and its output is shown verbatim
below it. Each file begins with \scm{(include "prelude.scm")}, which loads this prelude (the include line is
left out of the listings):
\begin{scheme}
(import scheme (chicken base) (chicken pretty-print)
        (aux base) (aux stream) (aux kanren micro))
(define-syntax-rule (show label expr)
  (begin (display label) (display " ") (write expr) (newline)))
\end{scheme}
\scm{(μkanren-run (q n grounded) g ...)} returns up to \scm{n} answers for \scm{q} (all answers
when \scm{n} is negative). With \scm{grounded} set to \scm{#t} it returns them as plain data; with
\scm{#f} it returns them as reified Scheme $\lambda$-expressions.
\end{tutorial}

% ======================================================================================================
\section{Which formulas are goals and which are clauses}\label{sec:hoprog-languages}

At first order (\cref{ch:horn,ch:hh}) goals and clauses of \fohh{} were
\begin{align*}
G &::= \top \mid A \mid G\wedge G \mid G\vee G \mid \forall x\,G \mid \exists x\,G \mid D\imp G\\
D &::= A \mid G\imp D \mid \forall x\,D \mid D\wedge D
\end{align*}
and \fohc{} is the fragment without $\forall$ and $\imp$ in goals \book{118}. The higher-order languages
keep this top-level shape. What changes is \emph{what an atom $A$ may be}: once quantification over
predicates is allowed, an atom may contain $\lambda$-terms and predicate variables, and its head may itself
be a variable.

\subsection{Rigid and flexible atoms}

\begin{definition}[Rigid and flexible atoms {\book{119}}]\label{def:hoprog-rigid}
The $\lambda$-normal form of a higher-order atomic formula has the shape $(h\,t_1\ldots t_n)$ where $h$ is
either a variable or a non-logical constant. The atom is \emph{rigid} if $h$ is a constant and
\emph{flexible} if $h$ is a variable.
\end{definition}

So \lp{age bob 23} and \lp{mappred P L K} are rigid, while \lp{P X Y} is flexible.
The atom \lp{(x\ age x 23) bob} is rigid too, because its $\lambda$-normal form is \lp{age bob 23}.
Flexible atoms in \emph{goals} are what make higher-order programming possible: in the body of
\lp{mappred P (X::L) (Y::K) :- P X Y, mappred P L K.} the goal \lp{P X Y} is flexible.
Flexible atoms as \emph{clause heads} are a different matter.

\subsection{Why clause heads must be rigid}\label{sec:hoprog-rigidheads}

The book gives two reasons for rejecting flexible heads \book{119--120}.

\paragraph{Consistency.} A program is \emph{inconsistent} if every formula is provable from it. Every
\fohc{} and \fohh{} program is consistent. This matters in practice: otherwise a complete interpreter would
have to check the program for consistency before it could answer ``no''. With flexible heads,
inconsistent programs are easy to write:
\begin{lprolog}
P 5 :- q.
q.
\end{lprolog}
To prove any goal $B$, backchain on the first clause with $P := \abs{x}{B}$; what is left is \lp{q}, which
is provable. The single clause $\forall p\,p$ is an even shorter example \book{120}.

\paragraph{Modularity.} A clause $\forall x_1\ldots\forall x_m(A_1\wedge\cdots\wedge A_n\imp A_0)$ is read as
part of the definition of the relation named by the head of $A_0$. If that head is a variable, the clause
adds meaning to \emph{every} predicate. That defeats modular reading of code, and it also defeats
indexing clauses by predicate name.

Two natural uses are lost, and the book argues that losing them is acceptable. \emph{Leibniz
equality} (two terms are equal if they satisfy the same properties) would turn $x+0=x$ into
\begin{lprolog}
P (X + 0) :- P X.
P X :- P (X + 0).
\end{lprolog}
This is logically meaningful but operationally a disaster. A goal containing $n$ occurrences of, say,
\lp{5 + 0} admits $2^n$ abstractions for \lp{P}, and even a goal with no occurrence can always
backchain on the first clause, so the search loops \book{120}. \emph{Computed clauses}, as in
\lp{?- convert d P, P => g.}, put a variable where a clause is expected. There is no simple way to
guarantee statically that the value bound to \lp{P} will be a legal clause \book{121}.

\subsection{No implications inside terms}\label{sec:hoprog-noimp}

Logical symbols inside the arguments of atoms are useful (the symmetric closure
\lp{x\y\ adj x y ; adj y x} is a term containing $\vee$), but an \emph{implication} inside a term breaks
goal-directed search. The book's counterexample \book{121} uses constants $r:\typ{o}\to\typ{o}$ and
$s,t:\typ{o}$, and a variable $Q:\typ{o}\to\typ{o}\to\typ{o}$:
\[
\exists Q\,\big[\,\forall p\,\forall q\,[\,r\,(p\imp q)\imp r\,(Q\,p\,q)\,]\ \wedge\ Q\,(t\vee s)\,(s\vee t)\,\big].
\]
The first conjunct forces $Q := \abs{x}{\abs{y}{x\imp y}}$: after the $\forall$ and $\imp$ introductions
the goal $r\,(Q\,p\,q)$ must match the hypothesis $r\,(p\imp q)$. The second conjunct then becomes
$(t\vee s)\imp(s\vee t)$, which is provable in any reasonable logic, but \emph{not} by the goal-directed
rules of \cref{ch:horn}: the instance of $Q$ moved a disjunction from a goal position into a hypothesis,
where a left rule for $\vee$ would be needed.

\begin{definition}[Polarity {\book{121--122}}]\label{def:hoprog-polarity}
When implications are excluded from terms, an occurrence of a logical symbol is \emph{positive} if and
only if it is either within an atomic formula or is in the scope of only logical constants and to the
left of an even number of implications. All other occurrences are \emph{negative}.
\end{definition}

The point of the definition is an invariant: during cut-free proof search, positive occurrences stay
positive and negative ones stay negative. Goal-directed search relies on it. In the counterexample both
disjunctions start out positive, and instantiating $Q$ makes one of them negative. (The ban is a design
choice, not the only one possible: the bibliographic notes recall Felty's result that $\imp$ may occur
in terms if $\vee$ and $\exists$ are banned from formulas and terms instead \book{148}.)

\subsection{Herbrand universes, hohc and hohh}\label{sec:hoprog-hohc}

The quantifier rules $\exists R$ and $\forall L$ substitute terms. To keep goals goals and clauses clauses,
the terms that may be substituted have to be restricted.

\begin{definition}[Herbrand universes {\book{122}}]\label{def:hoprog-herbrand}
Let $\Sig$ contain at least the logical constants. $\mathcal{H}_1^\Sig$ is the set of all (possibly open)
$\lambda$-normal terms over $\Sig$ that contain neither $\forall$ nor $\imp$. $\mathcal{H}_2^\Sig$ is the
set of all (possibly open) $\lambda$-normal terms over $\Sig$ that do not contain $\imp$.
\end{definition}

Terms of $\mathcal{H}_1^\Sig$ may contain $\top,\wedge,\vee,\exists$, and the types of their subterms are
unrestricted. $\mathcal{H}_2^\Sig$ also allows $\forall$.

\begin{definition}[\hohc{}: higher-order Horn clauses {\book{122}}]\label{def:hoprog-hohc}
Let $A$ range over atoms in $\mathcal{H}_1^\Sig$ and $A_r$ over \emph{rigid} atoms in
$\mathcal{H}_1^\Sig$. Goals and clauses are the formulas whose $\lambda$-normal forms are
\begin{align*}
G &::= \top \mid A \mid G\wedge G \mid G\vee G \mid \exists_\tau x\,G\\
D &::= A_r \mid G\imp D \mid D\wedge D \mid \forall_\tau x\,D
\end{align*}
with no restriction on $\tau$. The $G$-formulas are exactly the formulas of $\mathcal{H}_1^\Sig$ of type $\typ{o}$.
\end{definition}

The operational semantics is the goal-directed calculus of \cref{ch:horn}, with three
changes \book{123}: formulas are taken in $\lambda$-normal form, \rulename{initial} works up to
$\alpha$-conversion, and the quantifier rules instantiate with $\lambda$-terms from the Herbrand universe
($\Vdash$ is the typing judgement for the instantiation term):
\begin{mathpar}
\inferrule{ }{\Sig;\Prog\xrightarrow{A'}A}\;\rulename{initial}\ (A'\equiv_\alpha A)
\and
\inferrule{\seq{\Sig}{\Prog}{B[t/x]} \\ \Sig;\emptyset\Vdash t:\tau \\ t\in\mathcal{H}_1^\Sig}{\seq{\Sig}{\Prog}{\exists_\tau x\,B}}\;\rulename{$\exists$R}
\and
\inferrule{\Sig;\Prog\xrightarrow{D[t/x]}A \\ \Sig;\emptyset\Vdash t:\tau \\ t\in\mathcal{H}_1^\Sig}{\Sig;\Prog\xrightarrow{\forall_\tau x\,D}A}\;\rulename{$\forall$L}
\end{mathpar}
Why is the restriction needed? Without it, the goal $\exists p\,p$ could be instantiated with
$\forall x\,(q\,x)$ for some $q\in\Sig$, which is not an \hohc{} goal. Substituting a term of
$\mathcal{H}_1^\Sig$ into a goal or a clause keeps it in its category. Intuitively, we never substitute
something that could not appear as the argument of an atom or in the structure of a goal \book{123}. In an
implementation, $t$ is not guessed: it is a logic variable that higher-order unification later fixes in
\rulename{initial}.

\begin{theorem}[Completeness of goal-directed search for \hohc{} {\book{123--124}, \book{148}}]\label{thm:hoprog-hohc}
Derivability of goals from program clauses in this calculus coincides exactly with provability in
classical higher-order logic. By the bibliographic notes, this holds for both the classical and the
intuitionistic versions of Church's Simple Theory of Types, without the axioms of extensionality,
infinity and choice \cite{NadathurMiller1990}.
\end{theorem}

Since $\forall R$ and ${\imp}R$ never fire in \hohc{}, the signature and the program stay fixed during a
computation \book{124}.

\begin{definition}[\hohh{}: higher-order hereditary Harrop formulas {\book{124}}]\label{def:hoprog-hohh}
Let $A$ range over atoms in $\mathcal{H}_2^\Sig$ and $A_r$ over rigid atoms in $\mathcal{H}_2^\Sig$:
\begin{align*}
G &::= \top \mid A \mid G\wedge G \mid G\vee G \mid \exists_\tau x\,G \mid D\imp G \mid \forall_\tau x\,G\\
D &::= A_r \mid G\imp D \mid D\wedge D \mid \forall_\tau x\,D
\end{align*}
The instances for $\exists R$ and $\forall L$ are terms of $\mathcal{H}_2^\Sig$, where $\Sig$ is the
signature \emph{current at that point of the derivation}.
\end{definition}

Goals are richer than $\mathcal{H}_2^\Sig$ (they may contain $\imp$; no term of $\mathcal{H}_2^\Sig$ does).
Derivability now coincides with provability in higher-order \emph{intuitionistic}
logic \cite{MillerNadathurPfenningScedrov1991} \book{124}. Unlike \hohc{}, $\forall R$
(\rulename{generic}) and ${\imp}R$ (\rulename{augment}) are active, so the signature and the program change
during computation, as in \cref{ch:hh}.

\subsection{\texorpdfstring{hohh$^+$}{hohh+}: hiding predicate names}\label{sec:hoprog-hohhplus}

\Cref{ch:hh} could limit the \emph{scope} of an auxiliary predicate's clauses with $\imp$, but not the
\emph{visibility of its name}. A context that already had clauses for \lp{rev} could corrupt a
\lp{reverse} that assumed its own \lp{rev}. One small liberalisation fixes this.

\begin{definition}[\hohh$^+$ {\book{125}}]\label{def:hoprog-hohhplus}
\hohh$^+$ is \hohh{} in which the head of an atomic program clause may also be a \emph{variable}, provided that
variable is bound by an \emph{essentially universal} quantifier that has the clause in its scope.
\end{definition}

In a goal, a bound variable is essentially universal if it is bound by a positive $\forall$, a negative
$\exists$, or a term-level $\lambda$. It is essentially existential otherwise. The definition is dualised
for clauses (\book{205}; see \cref{ch:lambdatree}). Essentially universal variables only ever become
eigenvariables. When search reaches the goal $\forall p\,(\ldots D\ldots\imp G)$, $\forall R$ has already
replaced $p$ by a \emph{new constant}, so by the time $D$ is added to the program it is a legal \hohh{} clause
with a rigid head. Hence \hohh$^+$ does not change computational behaviour; it only lets a program
introduce \emph{closed, hidden} predicate definitions. The book's two versions of \lp{reverse}
\book{126}:
\begin{lprolog}
type reverse   list A -> list A -> o.
reverse L K :-
  pi rev\
    (pi L\ rev nil L L) &
    (pi X\ pi L\ pi K\ pi M\ rev (X::L) K M :- rev L K (X::M))
    => rev L K nil.
\end{lprolog}
\begin{lprolog}
type reverse   list A -> list A -> o.
reverse L K :- pi rv\
  ( rv nil K &
    (pi X\ pi N\ pi M\ rv (X::N) M :- rv N (X::M)))
  => rv L nil.
\end{lprolog}

\begin{tutorial}[What the interpreter does with \texttt{reverse (1::2::3::nil) K}]
Using the first version \book{125}:
\begin{enumerate}
\item backchain on \lp{reverse}; the body is a \lp{pi} goal, so $\forall R$ generates a new predicate
  symbol, call it \lp{c}, and adds it to the signature;
\item the body is now an implication, so ${\imp}R$ adds the two clauses
\begin{lprolog}
pi L\ c nil L L.
pi X\ pi L\ pi K\ pi M\ c (X::L) K M :- c L K (X::M).
\end{lprolog}
\item it proves \lp{c (1::2::3::nil) K nil} by ordinary backchaining and obtains \lp{K = (3::2::1::nil)};
\item on leaving the goal, it removes \lp{c} and its clauses from the context.
\end{enumerate}
No clause outside \lp{reverse} can mention \lp{c}: the name did not exist when those clauses were written.
\end{tutorial}

\begin{pitfall}[Shadowing and shared variables in assumed clauses]
In the first version the inner \lp{pi L\} and \lp{pi K\} \emph{shadow} the outer \lp{L} and \lp{K}. In
the second version the assumed fact \lp{rv nil K} mentions the \emph{outer} logic variable \lp{K}, so the
hypothetical clause is not closed: \lp{K} is shared with the caller and is \emph{not} renamed at each use of
the clause.
\end{pitfall}

\begin{kanren}[Hidden predicates are lexical scope]
In a Scheme-hosted miniKanren a relation is a procedure, and Scheme already has a way to make a name
invisible: a local definition. The first \lp{reverse}, with \lp{rev} hidden, is one \scm{let}:
\begin{scheme}
;; Figure 5.1 with the hidden predicate rev as a lexically local relation
(define reverse°
  (let ()
    (define-relation (rev l k m)
      (cond° ((null° l) (=° k m))
             ((fresh° (x l*) (cons° x l* l) (rev l* k (cons x m))))))
    (λ (l k) (rev l k '()))))
(show 1 (μkanren-run (q -1 #t) (reverse° '(1 2 3) q)))
(show 2 (μkanren-run (q 1 #t) (reverse° q '(1 2 3))))
\end{scheme}
\begin{answers}
1 ((3 2 1))
2 ((3 2 1))
\end{answers}
No client of \scm{reverse°} can add clauses to \scm{rev}. This is the encapsulation that \hohh$^+$ provides,
and it costs nothing. (The second query asks for one answer: in that mode a request for all answers
does not terminate.) What lexical scope does \emph{not} give is a clause set that \emph{grows during the
search}, such as the memo table of \cref{ex:hoprog-fib}, or a hidden name that exists as a first-class
\emph{constant}, one that can be stored in terms and must not escape its scope. Both need new machinery
in the state: a dynamic clause store for $\imp$, and eigenvariables with a scope check for $\forall$.
See \cref{ch:design}.
\end{kanren}

% ======================================================================================================
\section{Higher-order relational programming}\label{sec:hoprog-hoprog}

\subsection{Mapping, filtering and quantifying over lists}

The basic higher-order relations of the book \book{127}:
\begin{lprolog}
type foreach, forsome   (A -> o) -> list A -> o.
type mappred            (A -> B -> o) -> list A -> list B -> o.
type sublist            (A -> o) -> list A -> list A -> o.

foreach P nil.
foreach P (X::L) :- P X, foreach P L.

forsome P (X::L) :- P X; forsome P L.

mappred P nil nil.
mappred P (X::L) (Y::K) :- P X Y, mappred P L K.

sublist P (X::L) (X::K) :- P X, sublist P L K.
sublist P (X::L) K :- sublist P L K.
sublist P nil nil.
\end{lprolog}
Read relationally \book{126}: \lp{mappred P L K} says the lists have equal length and are related pointwise
by \lp{P}; \lp{forsome P L} says some member satisfies \lp{P}; \lp{foreach P L} says all members do;
\lp{sublist P L K} says \lp{K} is an order-preserving sublist of \lp{L} whose members all satisfy \lp{P}.
Because of the clause order, the \emph{maximal} sublists come first. A small database:
\begin{lprolog}
kind name    type.
type bob, sue, ned   name.
type age     name -> int -> o.
type male, female    name -> o.

age bob 23 & age sue 24 & age ned 23.
male bob & female sue & male ned.
\end{lprolog}
\begin{example}[Queries over the database {\book{126--127}}]\label{ex:hoprog-mappred}
\begin{lprolog}
?- mappred age (ned::bob::sue::nil) L.
L = (23::23::24::nil)
?- mappred age L (23::24::nil).
L = (bob::sue::nil);
L = (ned::sue::nil)
?- sublist male (ned::bob::sue::nil) L.
L = ned::bob::nil;
L = ned::nil;
L = bob::nil;
L = nil;
no
?- forsome female (ned::bob::sue::nil).
solved
?- foreach female (ned::bob::sue::nil).
no
\end{lprolog}
The second query runs \lp{mappred} ``backwards'': \lp{age} is a relation, not a function, so the list of
names is computed from the list of ages.
\end{example}

\subsection{\texorpdfstring{$\lambda$}{Lambda}-terms of predicate type as arguments}

A predicate argument need not be a constant. \lp{x\ age jane x} has type \lp{int -> o} (it is
$\eta$-equivalent to \lp{age jane}), and \lp{n\ age n 24} has type \lp{name -> o} \book{129}. The book
reports:
\begin{lprolog}
?- (n\ age n 24) W.          % β-equal to ?- age W 24.
W = sue
?- mappred (x\y\ age x y) (ned::bob::sue::nil) L.
L = (23::23::24::nil)
?- mappred (x\y\ age y x) (23::24::nil) K.
K = (bob::sue::nil);
K = (ned::sue::nil)
?- foreach (x\ sigma y\ age x y) (ned::bob::sue::nil).
solved
?- foreach (x\ age x A) (ned::bob::sue::nil).
no
?- foreach (x\ age x A) (ned::bob::nil).
A = 23
\end{lprolog}

\begin{pitfall}[A logic variable inside a \texorpdfstring{$\lambda$}{lambda} is shared, not abstracted]
In \lp{x\ age x A}, the variable \lp{A} lies \emph{outside} the abstraction: the same \lp{A} is used by
every instance of the body. The last two queries therefore ask for a \emph{common} age. Compare
\lp{x\ sigma y\ age x y}, where each element gets its own \lp{y}.
\end{pitfall}

\begin{kanren}[Predicate variables with known values are free]
In a Scheme-hosted miniKanren, a predicate argument whose value is known when it is called is just a
procedure that returns a goal. \lp{x\y\ age y x} becomes \scm{(λ (x y) (age° y x))}, and
\lp{x\ sigma y\ age x y} becomes \scm{(λ (x) (fresh° (y) (age° x y)))}. A logic variable captured by the
Scheme \scm{λ} is shared by every call, exactly as in the pitfall above:
\begin{scheme}
;; Figure 5.4: the age database, one relation per predicate constant
(define-relation (age° n a)
  (cond° ((=° n 'bob) (=° a 23))
         ((=° n 'sue) (=° a 24))
         ((=° n 'ned) (=° a 23))))
(define-relation (male° n) (or° (=° n 'bob) (=° n 'ned)))
;; Figure 5.3: the argument P is a Scheme procedure returning a goal
(define-relation (mappred° P L K)
  (cond° ((null° L) (null° K))
         ((fresh° (X L* Y K*) (cons° X L* L) (cons° Y K* K)
            (P X Y) (mappred° P L* K*)))))
(define-relation (foreach° P L)
  (cond° ((null° L))
         ((fresh° (X L*) (cons° X L* L) (P X) (foreach° P L*)))))
(define-relation (sublist° P L K)
  (cond° ((fresh° (X L* K*) (cons° X L* L) (cons° X K* K)
            (P X) (sublist° P L* K*)))
         ((fresh° (X L*) (cons° X L* L) (sublist° P L* K)))
         ((null° L) (null° K))))

(show 1 (μkanren-run (q -1 #t) (mappred° age° '(ned bob sue) q)))
(show 2 (μkanren-run (q -1 #t) (mappred° age° q '(23 24))))
(show 3 (μkanren-run (q -1 #t) (sublist° male° '(ned bob sue) q)))
;; λ-terms of predicate type are Scheme lambdas
(show 4 (μkanren-run (q -1 #t) (mappred° (λ (x y) (age° y x)) '(23 24) q)))
(show 5 (μkanren-run (q -1 #t)
          (foreach° (λ (x) (fresh° (y) (age° x y))) '(ned bob sue))
          (=° q 'solved)))
(show 6 (μkanren-run (q -1 #t) (foreach° (λ (x) (age° x q)) '(ned bob sue))))
(show 7 (μkanren-run (q -1 #t) (foreach° (λ (x) (age° x q)) '(ned bob))))
\end{scheme}
\begin{answers}
1 ((23 23 24))
2 ((bob sue) (ned sue))
3 (() (bob) (ned) (ned bob))
4 ((bob sue) (ned sue))
5 (solved)
6 ()
7 (23)
\end{answers}
Every answer \emph{set} matches the book. The \emph{order} of the \lp{sublist} answers (line 3) does not,
because \scm{cond°} interleaves its branches while λProlog searches them depth-first in clause order.
Compare answer sets unless the search strategy is part of the specification.
\end{kanren}

\subsection{Closures, composition and folds}\label{sec:hoprog-closures}

Relations that take relations and return relations \book{127--128}:
\begin{lprolog}
type ref, sym, trans   (A -> A -> o) -> A -> A -> o.
ref   R X Y :- X = Y; R X Y.
sym   R X Y :- R X Y; R Y X.
trans R X Y :- R X Y.
trans R X Z :- R X Y, trans R Y Z.

kind node  type.
type a, b, c, d, e   node.
type adj   node -> node -> o.
adj a b & adj b c & adj b d & adj d c & adj c e.
\end{lprolog}
Both \lp{?- trans adj a d.} and \lp{?- sym adj b a.} succeed. Thanks to currying, the equivalence closure
is just the partial application \lp{(trans (sym (ref adj)))}. The non-recursive closures can also be
written as $\lambda$-terms, for example \lp{(x\y\ x = y ; adj x y)}. The book adds a warning
\book{128}. The transitive closure can be written \emph{declaratively} as a single term,
\begin{lprolog}
x\y\ pi p\
   (pi U\ pi V\ adj U V => p U V) =>
   (pi U\ pi V\ pi W\ adj U V => p V W => p U W) =>
   p x y
\end{lprolog}
but it contains $\imp$. It is therefore \emph{not} in $\mathcal{H}_2^\Sig$ and cannot be passed as an
argument (\cref{sec:hoprog-noimp}).

Union, composition and a left fold \book{130}:
\begin{lprolog}
type union    (A -> B -> o) -> (A -> B -> o) -> A -> B -> o.
type compose  (A -> B -> o) -> (B -> C -> o) -> A -> C -> o.
type foldl    (A -> B -> B -> o) -> list A -> B -> B -> o.

union   R S X Y :- R X Y; S X Y.
compose R S X Y :- R X Z, S Z Y.
foldl P nil X X.
foldl P (Z::L) X Y :- P Z X W, foldl P L W Y.
\end{lprolog}
If \lp{P : A -> B -> B -> o} is a family of binary relations on \lp{B} indexed by \lp{A}, then
\lp{(foldl P)} composes them one after another along a list. Reversing a list then means pushing every
element onto a stack and popping them all off \book{130--131}:
\begin{lprolog}
kind stack   type -> type.
type emp     stack A.
type stk     A -> stack A -> stack A.
type empty   stack A -> o.
type enter, remove   A -> stack A -> stack A -> o.
type reverse list A -> list A -> o.

empty emp.
enter X S (stk X S).
remove X (stk X S) S.
reverse L K :- compose (foldl enter L) (foldl remove K) emp emp.
\end{lprolog}

\begin{kanren}[Partial application is curried \texttt{define}]
λProlog is curried, while \scm{define-relation} has a fixed arity. A relation-returning relation is
written with CHICKEN's curried \scm{define}, \scm{(define ((sym° R) x y) ...)}; a recursive one defines a
local relation and returns it:
\begin{scheme}
(define-relation (adj° x y)
  (cond° ((=° x 'a) (=° y 'b)) ((=° x 'b) (=° y 'c)) ((=° x 'b) (=° y 'd))
         ((=° x 'd) (=° y 'c)) ((=° x 'c) (=° y 'e))))
;; curried closure builders: (sym° R) is itself a binary relation
(define ((ref° R) x y) (or° (=° x y) (R x y)))
(define ((sym° R) x y) (or° (R x y) (R y x)))
(define (trans° R)
  (define-relation (T x z) (cond° ((R x z)) ((fresh° (y) (R x y) (T y z)))))
  T)
(show 1 (μkanren-run (q -1 #t) ((trans° adj°) 'a 'd) (=° q 'solved)))
(show 2 (μkanren-run (q -1 #t) ((sym° adj°) 'b 'a) (=° q 'solved)))
(show 3 (μkanren-run (q -1 #t) ((trans° adj°) 'a q)))
(show 4 (μkanren-run (q 1 #t)
          ((trans° (sym° (ref° adj°))) 'e 'a) (=° q 'solved)))
\end{scheme}
\begin{answers}
1 (solved)
2 (solved)
3 (b c d e c e)
4 (solved)
\end{answers}
Line 3 shows that \lp{trans} is a relation, not a set: \scm{c} and \scm{e} come back once per path. Line 4
uses the equivalence closure built by partial application. Composition and the stack-based reverse:
\begin{scheme}
;; Figures 5.5 and 5.6: compose, foldl, and reverse through a stack
(define ((compose° R S) x y) (fresh° (z) (R x z) (S z y)))
(define-relation (foldl° P L X Y)
  (cond° ((null° L) (=° X Y))
         ((fresh° (Z L* W) (cons° Z L* L) (P Z X W) (foldl° P L* W Y)))))
(define-relation (enter° X S S*) (=° S* `(stk ,X ,S)))
(define-relation (remove° X S S*) (=° S `(stk ,X ,S*)))
(define-relation (reverse° L K)
  ((compose° (λ (s s*) (foldl° enter° L s s*))
             (λ (s s*) (foldl° remove° K s s*)))
   'emp 'emp))
(show 1 (μkanren-run (q -1 #t) (reverse° '(1 2 3) q)))
(show 2 (μkanren-run (q -1 #t) (foldl° enter° '(1 2 3) 'emp q)))
\end{scheme}
\begin{answers}
1 ((3 2 1))
2 ((stk 3 (stk 2 (stk 1 emp))))
\end{answers}
The partial application \lp{(foldl enter L)} becomes an explicit \scm{λ}. None of this needs anything
beyond Scheme closures.
\end{kanren}

\subsection{Binary clauses: continuation passing}\label{sec:hoprog-cps}

Predicate variables of type $\typ{o}$ are \emph{goal-valued}, so a program can pass its own continuation
\book{131}. For each predicate $p:\tau_1\to\cdots\to\tau_j\to\typ{o}$ introduce
$p':\tau_1\to\cdots\to\tau_j\to\typ{o}\to\typ{o}$, and write $A'$ for $(p'\,t_1\ldots t_j)$ when $A=(p\,t_1\ldots t_j)$.
Then rewrite every clause:
\begin{gather*}
\forall\bar x\,[A_1\wedge\dots\wedge A_n\imp A_0]\ (n>0)\ \rightsquigarrow\ \forall\bar x\,\forall K\,[\,(A'_1\,(\dots(A'_n\,K)\dots))\imp(A'_0\,K)\,]\\
\forall\bar x\,A_0\ \rightsquigarrow\ \forall\bar x\,\forall K\,[\,K\imp(A'_0\,K)\,]
\end{gather*}
with $K:\typ{o}$. Conjunction becomes explicit sequencing, and every clause body has at most one, possibly
complex, atom. For example
\begin{lprolog}
adj a b.
adj b c.
path X Y :- adj X Y.
path X Y :- adj X Z, path Z Y.
\end{lprolog}
becomes
\begin{lprolog}
adj' a b K :- K.
adj' b c K :- K.
path' X Y K :- adj' X Y K.
path' X Y K :- adj' X Z (path' Z Y K).
\end{lprolog}
and \lp{?- path X Y.} is provable if and only if \lp{?- path' X Y true.} is provable \book{132}.

\begin{kanren}[Goals are already first-class]
In μKanren a goal is a value, so the binary-clause form is an idiom, not a feature. The continuation is a
goal argument:
\begin{scheme}
;; Section 5.3: binary clauses, the continuation K is a goal
(define (adj/k° x y k)
  (cond° ((=° x 'a) (=° y 'b) k) ((=° x 'b) (=° y 'c) k)))
(define-relation (path/k° x y k)
  (cond° ((adj/k° x y k))
         ((fresh° (z) (adj/k° x z (path/k° z y k))))))
(show 1 (μkanren-run (q -1 #t)
          (fresh° (x y) (path/k° x y ✓°) (=° q (list x y)))))
(show 2 (μkanren-run (q -1 #t)
          (fresh° (x y) (path/k° x y (=° q (list x y))))))
\end{scheme}
\begin{answers}
1 ((a b) (b c) (a c))
2 ((a b) (b c) (a c))
\end{answers}
In the second query the continuation does the work of the query itself. Conjunction in μKanren
(\scm{andª}, which feeds every state of the first goal to the second) is this same continuation
composition.
\end{kanren}

\subsection{Memoisation with hidden, growing predicates}

\begin{example}[Memoised Fibonacci in \hohh$^+$ {\book{132}}]\label{ex:hoprog-fib}
\begin{lprolog}
type fib_memo   int -> ((int -> int -> o) -> o) -> o.
fib_memo N G :-
   pi memo\ pi loop\
     memo 0 0 => memo 1 1 =>
     (pi F1\ pi F2\ pi F3\ pi C\ pi C'\
        (loop C F1 F2 :- (C >  N, (G memo)) ;
                         (C =< N, C' is C + 1, F3 is F1 + F2,
                          memo C F3 => loop C' F2 F3))) =>
     loop 2 0 1.
\end{lprolog}
\lp{memo} and \lp{loop} are hidden predicates, bound by \lp{pi} and used as heads of assumed clauses. The
memo table \emph{grows} through \lp{memo C F3 => ...}. When \lp{C > N}, the continuation \lp{G}, a
predicate that takes a predicate, is called with the hidden name \lp{memo}, whose extension is then
$\{(0,0),(1,1),\ldots,(N,\mathit{fib}\,N)\}$. For the query
\begin{lprolog}
?- fib_memo 20 (fib\ sigma M\ fib N M, M is N * N).
\end{lprolog}
the book states, without listing them, that ``there are exactly three values for N that make
this atom provable'' \book{132}. A short computation gives $N\in\{0,1,12\}$: the continuation asks
for $\mathit{fib}\,N=N\cdot N$, and $\mathit{fib}\,0=0$, $\mathit{fib}\,1=1$, $\mathit{fib}\,12=144$.
\end{example}

This example mixes three things: higher-order programming, dynamic clause addition, and a hidden name
handed to a continuation. A Scheme \scm{letrec} cannot express the growing table, which is why the
kanren box of \cref{sec:hoprog-hohhplus} lists dynamic clauses as a genuine requirement.

% ======================================================================================================
\section{Flexible goals}\label{sec:hoprog-flex}

What happens when a goal's head is a variable that is \emph{still unbound} when the goal is reached? Take
\lp{P bob 23} against the age database. \lp{P := age} works, but so do \lp{x\y\ age x 23, age ned y},
\lp{x\y\ age sue 24}, \lp{x\y\ memb 4 (3::4::5::nil)}, and in fact \lp{x\y\ g} for \emph{any} provable
closed goal \lp{g} \book{132--133}. Such a goal is \emph{underconstrained}. Three treatments are
possible \book{134}:
\begin{enumerate}
\item \textbf{Suspend} the goal in the hope that other goals will instantiate its head, and resume it when
  it becomes rigid. If only flexible goals remain, solve each of them by instantiating its head with
  $\abs{x_1}{\ldots\abs{x_n}{\top}}$, the universally true relation of the right type, which is ``the
  (extensionally) largest one that works''.
\item \textbf{Eagerly} bind the head to the universally true relation.
\item Issue a \textbf{run-time error}.
\end{enumerate}
The first option is complete with respect to provability in intuitionistic higher-order logic. The other
two are incomplete. Yet experience with λProlog implementations says the third is the most useful while
developing code: a flexible goal met by depth-first search is usually a bug, for example a predicate
constant mistyped with a capital letter \book{134}. Option~1 is the same idea that unification-based
interpreters apply to flexible-flexible unification problems: the completeness results ``provide the
foundations for a unification-based interpreter \ldots\ that explicitly delays the consideration of
flexible-flexible unification problems as well as flexible goals'' \book{148}.

A predicate variable becomes meaningful when the program first \emph{constrains} it to a structured space
of candidates \book{133}:
\begin{lprolog}
kind i  type.
type jane, mary, john   i.
type mother, father, wife, husband   i -> i -> o.
type primrel, rel   (i -> i -> o) -> o.

primrel father & primrel mother & primrel wife & primrel husband.
rel R :- primrel R.
rel (x\y\ sigma z\ R x z, S z y) :- primrel R, primrel S.
mother jane mary & wife john jane.
\end{lprolog}
\lp{?- R john mary.} on its own is underconstrained. \lp{?- rel R, R john mary.} is meaningful:
\lp{rel} enumerates candidate relations, after which \lp{R john mary} is rigid. The only answer is
\lp{R = x\y\ sigma Z\ wife x Z, mother Z y}. Note that the head of the second \lp{rel} clause contains a
$\lambda$-term with $\exists$ and $\wedge$. This is legal because the term lies in $\mathcal{H}_1^\Sig$.
Whether a flexible goal is underconstrained thus ``depends on when the goal is invoked''
\book{133}: goal order matters operationally.

\begin{tutorial}[Depth-first search for \texttt{rel R, R john mary}]
Rename the clause variables of the second \lp{rel} clause to \lp{R1} and \lp{S1}, so that they are not
confused with the query's \lp{R}.
\begin{enumerate}
\item \lp{rel R} backchains on the first clause and \lp{primrel R} binds \lp{R := father}. The second
  goal is now the \emph{rigid} atom \lp{father john mary}, which has no clauses and fails. Backtracking
  tries \lp{mother}, \lp{wife} and \lp{husband} in turn; all fail the same way.
\item \lp{rel R} backchains on the second clause:
  \lp{R := x\y\ sigma z\ R1 x z, S1 z y}, with the subgoals \lp{primrel R1} and \lp{primrel S1}.
  After they bind \lp{R1} and \lp{S1}, the goal \lp{R john mary} is $\beta$-reduced to
  \lp{sigma z\ R1 john z, S1 z mary}, which is again built from rigid atoms.
\item With \lp{R1 := father} or \lp{R1 := mother} the atom \lp{R1 john z} fails for every choice of
  \lp{S1}. With \lp{R1 := wife} it gives \lp{z := jane}; then \lp{S1 := father} fails on
  \lp{father jane mary} and \lp{S1 := mother} succeeds on \lp{mother jane mary}.
\end{enumerate}
The answer is \lp{R} with \lp{R1} and \lp{S1} instantiated, which is the book's answer. The predicate
variable is never invoked while unbound: \lp{rel} always gives it a value first.
\end{tutorial}

\begin{kanren}[Binding a logic variable to a relation: possible, but opaque]
Any Scheme value can be bound to a logic variable, procedures included, so \lp{rel} can be transcribed.
The composed relation is built by \scm{project°} once \scm{r} and \scm{s} are known, and invoking the
value of \scm{R} means walking it and applying it:
\begin{scheme}
(import (chicken condition))
;; Figure 5.8: relations as values of logic variables
(define-relation (mother° x y) (=° x 'jane) (=° y 'mary))
(define-relation (wife° x y) (=° x 'john) (=° y 'jane))
(define-relation (father° x y) ✗°)
(define-relation (husband° x y) ✗°)
(define-relation (primrel° R)
  (cond° ((=° R father°)) ((=° R mother°)) ((=° R wife°)) ((=° R husband°))))
(define-relation (rel° R)
  (cond° ((primrel° R))
         ((fresh° (r s) (primrel° r) (primrel° s)
            (project° ((r* r) (s* s))
              (=° R (λ (x y) (fresh° (z) (r* x z) (s* z y)))))))))
;; invoking the value of R: walk it, then apply it
(define-relation (rel-john-mary° R)
  (rel° R) (project° ((R* R)) (R* 'john 'mary)))
(show 1 (μkanren-run (q -1 #f) (rel-john-mary° q)))
(show 2 (μkanren-run (q -1 #f) (fresh° (R) (rel-john-mary° R) (=° q 'found))))
;; a flexible goal: R is still unbound when it is invoked
(show 3 (handle-exceptions exn
            (list 'error ((condition-property-accessor 'exn 'message) exn))
          (μkanren-run (q -1 #t) (project° ((R* q)) (R* 'bob 23)))))
\end{scheme}
\begin{answers}
1 ((λ () #<procedure (? x y)>))
2 ((λ () (quote found)))
3 (error "call of non-procedure")
\end{answers}
The search finds the single answer (line 2), but the answer itself is \scm{#<procedure (? x y)>}, not
\lp{x\y\ sigma Z\ wife x Z, mother Z y}. Grounded reification (\scm{#t}) even fails on it, because it
\scm{eval}s the reified expression, and \scm{eval} rejects a procedure as an ``illegal non-atomic
object''. Line 3 is a flexible goal: \scm{project°} hands over an unbound variable record, and applying
it is a Scheme error. That is the book's option~3, but only by accident. Options~1 and~2 need a goal that
can be \emph{inspected} before it is run. Option~1 also needs a store of suspended goals that unification
wakes up, much as the disequality store is re-checked after each \scm{=°}. Printing \lp{R} as a
$\lambda$-term needs relations represented as terms. The proposal keeps relations as closures
(\cref{sec:roadmap-lim-dynamic}); predicate quantification is sketched as \cref{sec:roadmap-later-pred}.
\end{kanren}

% ======================================================================================================
\section{Reasoning about programs by substitution}\label{sec:hoprog-reason}

Logic gives a means of \emph{reasoning} about programs, and hidden predicates make it sharper
\book{134--136}. Claim: the second \lp{reverse} of \cref{sec:hoprog-hohhplus} is symmetric, that is, if
\lp{reverse L K} is provable then so is \lp{reverse K L}. The proof outline \book{135--136}:
\begin{enumerate}
\item The only way to prove \lp{reverse L K} is by backchaining, so its body is provable:
\begin{lprolog}
pi rv\ ( rv nil K &
         (pi X\ pi N\ pi M\ rv (X::N) M :- rv N (X::M)))
       => rv L nil
\end{lprolog}
\item A provable universal has provable instances. Instantiate \lp{rv} with \lp{x\y\ (not (rv y x))}, where
  \lp{not} is \emph{logical} negation, not negation as failure:
\begin{lprolog}
not(rv K nil) &
(pi X\ pi N\ pi M\ not(rv M (X::N)) :- not(rv (X::M) N))
=> not(rv nil L).
\end{lprolog}
\item Take the contrapositive of each implication, which is classically valid:
\begin{lprolog}
rv nil L & (pi X\ pi N\ pi M\ rv (X::M) N :- rv M (X::N))
=> rv K nil.
\end{lprolog}
\item This is a first-order Horn entailment, and on Horn clauses classical and intuitionistic provability
  coincide.
\item Generalise over \lp{rv} and rename \lp{M} and \lp{N}. The result is the body of \lp{reverse} with
  \lp{L} and \lp{K} swapped. \qed
\end{enumerate}
Intuitively, the trace of \lp{rv} on \lp{(a::b::c::nil)} is a table with rows
\lp{(a::b::c::nil) nil}, \lp{(b::c::nil) (a::nil)}, \lp{(c::nil) (b::a::nil)}, \lp{nil (c::b::a::nil)}.
Swapping the columns and reading the rows in reverse order gives a valid trace for the reverse of
\lp{(c::b::a::nil)} \book{135}.

Three subtleties are worth noting. The \lp{rv} inside the instantiating term \lp{x\y\ (not (rv y x))}
is a \emph{new} predicate variable that happens to reuse the name; step~5 generalises over it. The
instantiating term contains $\neg$, that is $\imp\bot$, so it is
\emph{not} in $\mathcal{H}_2^\Sig$. That is fine, because the argument uses the metatheory
(soundness, then instantiation, contraposition and completeness), not goal-directed search. And the argument
needs \lp{rv} to be hidden: no other clause can mention it, so any $\lambda$-term may be ``plugged in''
at its place.

\begin{kanren}[Not available to closures]
A relation that is a Scheme closure offers nothing to substitute into: the argument above treats the
program as a \emph{formula}. A port can at most turn the theorem into a property test (for random
lists, \lp{reverse L K} implies \lp{reverse K L}). Reasoning of this kind belongs to a deep
embedding, where clauses are data.
\end{kanren}

% ======================================================================================================
\section{Defining the logical constants}\label{sec:hoprog-logconst}

Predicate variables are powerful enough to \emph{define} logical constants inside the logic
\book{136--137}. In \hohc{}, clauses can specify the right-introduction rules:
\begin{lprolog}
type tt, ff   o.
type or       o -> o -> o.
type exists   (A -> o) -> o.
tt.                  % true
or P Q :- P.         % disjunction
or P Q :- Q.
exists B :- B T.     % existential quantifier
\end{lprolog}
\lp{ff} has no clauses, so proving it fails. Only the right rules are captured. The left rules (for
example the rule of cases for a disjunctive hypothesis) are not. In \hohh{} and \hohh$^+$, $\forall$
and $\imp$ give the standard impredicative encodings \book{136}:
\begin{gather*}
\bot \equiv \forall p\,p \qquad
\top \equiv \forall p\,(p\imp p)\\
B\vee C \equiv \forall p\,((B\imp p)\imp(C\imp p)\imp p)\qquad
\exists x\,(G\,x) \equiv \forall p\,((\forall x\,(G\,x\imp p))\imp p)
\end{gather*}
(the last three are \hohh$^+$ goals, since $p$ heads assumed clauses and is bound by a positive $\forall$).

\begin{tutorial}[Running the encodings]
Goal-directed search reads them as follows. For $\forall p\,p$, $p$ becomes a fresh constant with no
clauses, so the goal fails. For $\forall p\,((B\imp p)\imp(C\imp p)\imp p)$, a fresh $c$ is created, the
clauses $c$ \lp{:-} $B$ and $c$ \lp{:-} $C$ are assumed, and $c$ is proved by backchaining on either clause,
that is, by proving $B$ or $C$. For $\forall p\,((\forall x\,(G\,x\imp p))\imp p)$, a fresh $c$ is created,
the clause $\forall x\,(c$ \lp{:-} $G\,x)$ is assumed, and $c$ is proved by backchaining: $x$ becomes a logic
variable $X$ and the goal $G\,X$ remains.
\end{tutorial}

\begin{kanren}[Which definitions are free]
The \hohc{} definitions are one-liners over existing goals: \lp{tt} is \scm{✓°}, \lp{ff} is \scm{✗°},
\lp{or} is \scm{or°}, and \lp{exists B :- B T} is \scm{(fresh° (t) (B t))}. The \hohh{} encodings are not
free. They need a goal-level $\forall$ that creates a fresh constant and a goal-level $\imp$ that assumes
clauses for it. Those are the two state extensions of the kanren box in \cref{sec:hoprog-hohhplus}.
\end{kanren}

% ======================================================================================================
\section{Conditional and negation-as-failure}\label{sec:hoprog-naf}

With the cut \lp{!}, the book defines a conditional and negation-as-failure \book{137}:
\begin{lprolog}
type if   o -> o -> o -> o.
if P Q R :- P, !, Q.
if P Q R :- R.

type not  o -> o.
not P :- P, !, fail.
not P.
% alternatively
not P :- if P fail true.
\end{lprolog}
\lp{if P Q R} tries \lp{P}. If it succeeds, the cut commits to its first solution and \lp{Q} is proved;
only if \lp{P} fails is \lp{R} proved. This is not logic: \lp{?- X = 2, not (1 = X).} succeeds, while
\lp{?- not (1 = X), X = 2.} fails. ``Thus even a basic property such as commutativity of conjunctions
becomes invalid in this setting'', and the book uses these predicates ``extremely sparingly''
\book{137--138}.

\begin{kanren}[The cut as \texttt{take° 1}]
\micro{} has the soft cut \scm{if°}, which keeps \emph{every} answer of its test (like miniKanren's
\texttt{conda}), and \scm{take°}. The cut-based \lp{if} commits to the first answer, so the faithful
transcription of \lp{not P} is \scm{(if° (take° 1 P) ✗° ✓°)}. The \scm{take° 1} is not cosmetic: without it,
a test with infinitely many answers would make the conditional fail on each of them forever. The
non-commutativity reproduces exactly, and the disequality constraint \scm{≠°} is the declarative
alternative that commutes:
\begin{scheme}
;; Section 5.7: not P :- if P fail true, the cut read as "first answer only"
(show 1 (μkanren-run (q -1 #t) (=° q 2) (if° (take° 1 (=° 1 q)) ✗° ✓°)))
(show 2 (μkanren-run (q -1 #t) (if° (take° 1 (=° 1 q)) ✗° ✓°) (=° q 2)))
;; the declarative alternative: a disequality constraint commutes
(show 3 (μkanren-run (q -1 #t) (=° q 2) (≠° 1 q)))
(show 4 (μkanren-run (q -1 #t) (≠° 1 q) (=° q 2)))
\end{scheme}
\begin{answers}
1 (2)
2 ()
3 (2)
4 (2)
\end{answers}
\end{kanren}

% ======================================================================================================
\section{\texorpdfstring{$\lambda$}{Lambda}-terms as functions}\label{sec:hoprog-functions}

So far the higher-order variables have been predicate variables. Variables of non-predicate functional
type are instantiated by $\lambda$-terms that are only \emph{evaluated}, by $\beta\eta$-conversion. As
\cref{ch:lambda} showed, these functions are weak. ``This weakness, however can be an asset for a
different kind of computation: Rather than instantiating function variables with $\lambda$-terms, we can
think of finding values for them by solving unification problems'' \book{138}.

\subsection{\texttt{mapfun} and \texttt{reducefun}}

\begin{lprolog}
type mapfun     (A -> B) -> list A -> list B -> o.
type reducefun  (A -> B -> B) -> list A -> B -> B -> o.

mapfun F nil nil.
mapfun F (X::L) ((F X)::K) :- mapfun F L K.

reducefun F nil Z Z.
reducefun F (H::T) Z (F H R) :- reducefun F T Z R.
\end{lprolog}
Assume a type \lp{i}, constants \lp{a, b, c, d : i} and \lp{g : i -> i -> i}.

\paragraph{Forwards.} \lp{?- mapfun (x\ g a x) (a::b::nil) L.} has exactly one answer,
\lp{L = ((g a a)::(g a b)::nil)}, obtained by $\beta$-normalising \lp{((x\ g a x) a)} and
\lp{((x\ g a x) b)} \book{138}. \lp{mappred} runs arbitrary goals, so it is strictly more powerful. In fact
\lp{mapfun F L K :- mappred (x\y\ y = F x) L K.}

\paragraph{Backwards: function synthesis.} \book{139}
\begin{lprolog}
?- mapfun F (a::b::nil) ((g a a)::(g a b)::nil).
F = x\ g a x                        % the only answer
\end{lprolog}
The first element alone poses \lp{F a = g a a}, which has four unifiers:
\begin{lprolog}
(x\ g x x)   (x\ g a x)   (x\ g x a)   (x\ g a a)
\end{lprolog}
Only \lp{x\ g a x} also solves \lp{F b = g a b}, so an interpreter that tries another unifier first has to
backtrack. Weakness can also show up as failure: \lp{?- mapfun F (a::b::nil) (c::d::nil).} answers
\lp{no}. Simply typed $\lambda$-terms cannot do case analysis on constants, so no term maps
\lp{a} to \lp{c} and \lp{b} to \lp{d} (\cref{ch:lambda}) \book{139--140}.

\paragraph{\texttt{reducefun}.} Here \lp{+} is only a term constructor, with no arithmetic \book{140}:
\begin{lprolog}
?- reducefun (x\y\ x + y) (3::4::8::nil) 6 R.
R = 3 + (4 + (8 + 6));
?- reducefun F (4::8::nil) 6 (1 + (4 + (1 + (8 + 6)))).
F = x\y\ 1 + (4 + (1 + (8 + 6)));
F = x\y\ 1 + (x + (1 + (8 + 6)));
F = x\y\ 1 + (x + y);
no
?- pi z\ reducefun F (4::8::nil) z (1 + (4 + (1 + (8 + z)))).
F = x\y\ 1 + (x + y);
no
\end{lprolog}

\begin{tutorial}[Why three answers, and then one]
Unfolding the two clauses turns the second query into the single equation
$F\,4\,(F\,8\,6) = 1+(4+(1+(8+6)))$. Its solutions are the bodies obtained by abstracting some occurrences
of the constants of the target consistently. The constant function abstracts nothing. The second answer
abstracts the \lp{4}. The third abstracts both the \lp{4} and the inner \lp{1 + (8 + 6)}, which has the
shape \lp{x + y}. A candidate such as \lp{x\y\ 1 + (4 + y)} fails, because $F\,8\,6$ would give
$1+(4+6)\neq 1+(8+6)$. In the third query, $6$ is replaced by a new constant $z$ introduced by
\lp{pi} \emph{inside} the scope of $F$'s implicit $\exists$. Instantiations of $F$ may not contain $z$, and
that eigenvariable check removes the two answers that would have to copy the base case into $F$. Only the
real fold function is left.
\end{tutorial}

\begin{kanren}[Forwards is free, synthesising \texttt{F} is not]
With \lp{F} a Scheme procedure that \emph{builds a term}, \lp{mapfun} and \lp{reducefun} run
forwards unchanged. Since \lp{F} is applied to a logic variable while the goal is being built, they also
run backwards in the \emph{list} argument:
\begin{scheme}
;; Figure 5.10, forward mode only: F is a Scheme procedure that builds a term
(define-relation (mapfun° F L K)
  (cond° ((null° L) (null° K))
         ((fresh° (X L* K*) (cons° X L* L) (cons° (F X) K* K)
            (mapfun° F L* K*)))))
(define-relation (reducefun° F L Z R)
  (cond° ((null° L) (=° Z R))
         ((fresh° (H T R*) (cons° H T L) (=° R (F H R*))
            (reducefun° F T Z R*)))))
(show 1 (μkanren-run (q -1 #t) (mapfun° (λ (x) `(g a ,x)) '(a b) q)))
(show 2 (μkanren-run (q -1 #t)
          (reducefun° (λ (x y) `(+ ,x ,y)) '(3 4 8) 6 q)))
;; F is applied to a logic variable while the goal is built:
;; the term (g a X) contains X, so the list argument can be unknown
(show 3 (μkanren-run (q -1 #t)
          (mapfun° (λ (x) `(g a ,x)) q '((g a a) (g a b)))))
\end{scheme}
\begin{answers}
1 (((g a a) (g a b)))
2 ((+ 3 (+ 4 (+ 8 6))))
3 ((a b))
\end{answers}
What cannot be transcribed is the case where \lp{F} itself is unknown. A Scheme procedure is opaque, so
there is nothing to unify \lp{F a} against. Function synthesis needs $\lambda$-terms as data, with
$\beta$-normalisation, and a unifier that handles applications of a variable. The problems above,
\lp{F a = g a a} and \lp{F 4 R = ...}, are \emph{not} higher-order patterns in the sense of
\cite{Miller1991}, because \lp{F} is applied to constants. The first is a second-order
\emph{matching} problem (its right-hand side is closed); the second is close to one (the flexible side
also contains the unknown \lp{R}). Both need Huet-style branching \cite{Huet1975} (see \cref{ch:hounif}). \Cref{sec:roadmap-later-huet} discusses where that
belongs in \micro{}.
\end{kanren}

\subsection{Functional difference lists}\label{sec:hoprog-fdl}

A first-order difference list \lp{(dl L K)}, with \lp{K} a suffix of \lp{L}, denotes \lp{L} without
\lp{K}. For instance, both of the terms
\lp{(dl (1::2::3::4::5::nil) (4::5::nil))} and \lp{(dl (1::2::3::L) L)}
denote the list \lp{(1::2::3::nil)}, and with the second form concatenation needs no recursion. A \emph{functional}
difference list (due to Brisset and Ridoux \cite{BrissetRidoux1991}) represents the list as a function
from its tail to the whole list \book{140--142}. The book's Figure~5.11 collects the labels of a
binary tree into a list in two ways, with ordinary and with functional difference lists. Here are the
declarations and the functional version:
\begin{lprolog}
kind dlist    type -> type.
type dl       list A -> list A -> dlist A.

type concat   dlist A -> dlist A -> dlist A -> o.
concat (dl L1 L2) (dl L2 L3) (dl L1 L3).

kind btree    type -> type.
type empty    btree A.
type bt       A -> btree A -> btree A -> btree A.

kind fdlist   type -> type.
type fdl      (list A -> list A) -> fdlist A.

type collect'  btree A -> list A -> o.
type aux'      btree A -> fdlist A -> o.
collect' Bt (A nil) :- aux' Bt (fdl A).
aux' empty (fdl x\ x).
aux' (bt N L R) (fdl x\ A (N::(B x))) :- aux' L (fdl A),
                                         aux' R (fdl B).
\end{lprolog}
\lp{(fdl x\ (1::2::3::x))} denotes \lp{(1::2::3::nil)}, with a ``pointer'' to its tail. Each operation is
a single unification \book{142}:
\begin{itemize}
\item head and tail: unify with \lp{x\ (H::(T x))}; the head is \lp{H} and the tail is \lp{(fdl T)};
\item last and front: unify with \lp{x\ F (Y::x)}; the last element is \lp{Y} and the front is \lp{(fdl F)};
\item concatenation of \lp{(fdl A)} and \lp{(fdl B)}: \lp{(fdl x\ A (B x))}, that is, function composition;
\item conversion to a list: apply to \lp{nil}, as \lp{(A nil)} in \lp{collect'}.
\end{itemize}
Palindromes then need no explicit reversal \book{142}:
\begin{lprolog}
type palindrome   fdlist A -> o.
palindrome (fdl x\x).
palindrome (fdl x\ Y::x).
palindrome (fdl x\ Y::(F (Y::x))) :- palindrome (fdl F).
\end{lprolog}
\begin{lprolog}
?- palindrome (fdl x\ a::b::c::b::a::x).
solved
?- palindrome (fdl x\ a::b::a::(F x)).
F = x\ x;
F = x\ b::a::x;
F = x\ a::b::a::x;
F = x\ Y::a::b::a::x;
F = x\ Z::Z::a::b::a::x;
F = x\ Y::Z::Y::a::b::a::x;
F = x\ Y::Z::Z::Y::a::b::a::x;
F = x\ Y::Z::U::Z::Y::a::b::a::x;
F = x\ Y::Z::U::U::Z::Y::a::b::a::x
\end{lprolog}
The second query has infinitely many answers. It finds the \lp{F} whose list, appended after
\lp{a::b::a}, makes a palindrome.

\begin{pitfall}[Types do not enforce the invariants]
\lp{(dl (1::2::3::nil) (4::5::nil))} is well typed, but it does not denote a list in the intended sense.
Likewise a function that ignores its argument, as in \lp{(fdl x\ (1::2::3::nil))}, has no tail ``pointer''
\book{143}. (The book prints the second term as \lp{(fdl x\ (1::2::3))}, which is literally ill typed.)
\end{pitfall}

\begin{kanren}[A closure is an fdl you can only run]
In its intended mode (tree given, list wanted), \lp{collect'} needs no higher-order unification. Every
clause head is unified with a fresh variable, so the only requirement is $\lambda$-terms that can be
built and applied. A Scheme closure provides that, using \scm{project°} to compose once the parts are
known:
\begin{scheme}
;; Figure 5.11, collect' in its intended mode, an fdl as a Scheme closure
(define-relation (aux/fdl° bt F)
  (cond° ((=° bt 'empty) (=° F (λ (x) x)))
         ((fresh° (n l r a b) (=° bt `(bt ,n ,l ,r))
            (aux/fdl° l a) (aux/fdl° r b)
            (project° ((a* a) (b* b))
              (=° F (λ (x) (a* (cons n (b* x))))))))))
(define-relation (collect/fdl° bt L)
  (fresh° (a) (aux/fdl° bt a) (project° ((a* a)) (=° L (a* '())))))
(define tree '(bt 2 (bt 1 empty empty) (bt 3 empty empty)))
(show 1 (μkanren-run (q -1 #t) (collect/fdl° tree q)))
\end{scheme}
\begin{answers}
1 ((1 2 3))
\end{answers}
\scm{project°} makes this non-relational: it only works when the subtrees are known. The operations that
make functional difference lists interesting (head and tail by unification with \lp{x\ (H::(T x))}, last
and front with \lp{x\ F (Y::x)}, \lp{palindrome}, and \lp{collect'} run backwards) all unify a pattern
\emph{against} a $\lambda$-term. They need $\lambda$-terms in the term language, and \lp{x\ F (Y::x)}
already applies a variable to a non-variable argument (\cref{ch:hounif}).
\end{kanren}

% ======================================================================================================
\section{Higher-order unification is not a panacea}\label{sec:hoprog-panacea}

It is tempting to make higher-order unification do all the work. The book's examples show why that is
often a mistake \book{143--145}. Take \lp{kind i type. type a,b i. type f i -> i -> i.}. The aim is to
abstract all occurrences of \lp{a} in \lp{(f a (f a b))}, that is, to obtain \lp{(x\ f x (f x b))}. The
naive program:
\begin{lprolog}
type extract_a  i -> (i -> i) -> o.
extract_a (F a) F.

?- extract_a (f a (f a b)) F.
F = x\ f x (f x b);
F = x\ f x (f a b);
F = x\ f a (f x b);
F = x\ f a (f a b);
no
\end{lprolog}
With $n$ occurrences there are $2^n$ ``extractions''. A cut, \lp{extract_a (F a) F :- !.}, relies on the
order in which the unifiers are enumerated, so it is not declarative. Two better options:
\begin{itemize}
\item \textbf{Reorder the quantifiers} so that the eigenvariable check does the selection:
\begin{lprolog}
pi a\ sigma F\ (F a) = (f a (f a b)).   % four proofs
sigma F\ pi a\ (F a) = (f a (f a b)).   % exactly one: F = x\ (f x (f x b))
\end{lprolog}
In the second line \lp{a} is a constant that is new for \lp{F}, so \lp{F} cannot mention it and every
occurrence must be abstracted. (This is also a higher-order pattern, \lp{F} applied to a distinct bound
variable, which is what makes it well behaved; see \cref{ch:lambdatree}.)
\item \textbf{Recurse on the syntax}, using the signature of \lp{i}:
\begin{lprolog}
extract_a a x\x.
extract_a b x\b.
extract_a (f T S) (x\ f (U x) (V x)) :- extract_a T U, extract_a S V.
\end{lprolog}
\end{itemize}
A second example is rewriting by congruence \book{145}:
\begin{lprolog}
type rewrite  int -> int -> o.
rewrite (0 + X) X.
rewrite (1 * X) X.
rewrite (X - X) 0.
rewrite (C X) (C Y) :- rewrite X Y.
\end{lprolog}
The last clause uses \lp{C} as a ``context with a hole''. \lp{rewrite ((1 * 5) - 5) (5 - 5)} is derivable,
but it has \emph{infinitely many} derivations: the last clause can be used any number of times with
\lp{C = x\x} before \lp{C = x\ x - 5} is finally chosen. The remedy is again explicit recursion over the
structure of terms.

\begin{kanren}[Interleaving as a home for non-unitary unification]
Every example in this section is a case where unification returns a \emph{stream} of solutions. In a
miniKanren, alternative unifiers are naturally branches of a disjunction, and interleaving enumerates
infinite families fairly. That fits Huet's procedure \cite{Huet1975} well. The lesson of the book still
applies: design relations so that unification has at most one answer (patterns, structural recursion, the
right quantifier order), and treat branching unification as a last resort. \Cref{ch:hounif} and
\cref{sec:roadmap-later-huet} take this up.
\end{kanren}

% ======================================================================================================
\section{Comparison with functional programming}\label{sec:hoprog-fp}

Most higher-order idioms of Scheme or ML carry over to logic programming easily (\lp{mappred},
\lp{sublist}, folds). The book argues that the logic programming notion is stronger in two ways
\book{146--147}:
\begin{enumerate}
\item predicate variables are only one kind of functional variable. Non-predicate function variables
  give new kinds of computation (\cref{sec:hoprog-functions}), developed in \cref{ch:lambdatree} for
  syntax with binders;
\item predicate expressions can be \emph{compared}, intensionally:
\begin{lprolog}
type eq_pred (A -> o) -> (A -> o) -> o.
eq_pred R R.
\end{lprolog}
  succeeds if and only if its arguments are unifiable predicate \emph{expressions}. For example,
  \lp{?- eq_pred (x\ p x, q x) (x\ q x, p x).} fails whatever the clauses for \lp{p} and \lp{q}. The
  comparison is structural, modulo $\lambda$-conversion, not equality of sets. Functional languages forbid
  comparing functions because their semantics is extensional, and extensional equality is undecidable.
\end{enumerate}
A design remark follows \book{147}. If the only goal were functional-style higher-order programming, a
more restricted syntax would do, for example one that forbids quantification over higher-order
non-predicate types. What general higher-order quantification adds is $\lambda$-terms that represent
syntax with binding, together with functional quantification to manipulate them. That was ``much of the
impetus for the development of λProlog''. Nor is λProlog a merger of functional and logic
programming: typed $\lambda$-terms are too weak as functions for that \book{148}.

\begin{kanren}[Unification on closures is pointer identity]
\micro{} unifies procedures by \scm{eq?}, which is finer than λProlog's $\lambda$-convertibility:
\begin{scheme}
;; Section 5.10: unification on procedures is pointer identity (eq?)
(define-relation (age° n a)
  (cond° ((=° n 'bob) (=° a 23)) ((=° n 'sue) (=° a 24))))
(define (older-than n)
  (λ (x) (fresh° (a) (age° x a) (project° ((a* a)) (if (> a* n) ✓° ✗°)))))
(show 1 (μkanren-run (q -1 #t) (=° age° age°) (=° q 'same)))
(show 2 (μkanren-run (q -1 #t)
          (=° (older-than 20) (older-than 20)) (=° q 'same)))
(show 3 (μkanren-run (q -1 #t)
          (let ((p (older-than 20))) (=° p p)) (=° q 'same)))
\end{scheme}
\begin{answers}
1 (same)
2 ()
3 (same)
\end{answers}
Two evaluations of the same \scm{λ} with a captured value give distinct closures (line 2), whereas
\lp{eq_pred (x\ p x) (x\ p x)} succeeds in λProlog. Intensional comparison, and every feature that
inspects a predicate (printing it, matching on it, \lp{eq_pred}), needs predicates represented as terms.
\end{kanren}

% ======================================================================================================
\section{Modules: programming in the large with logic}\label{sec:hoprog-modules}

Languages usually start with programming-in-the-small, and later ``a second modules language is often
imposed on the previously existing core language''. Its constructs may have ``little or no connection''
with the semantics of the core \book{150}. Chapter~6 of the book argues the opposite for logic
programming: the abstractions needed for modularity are \emph{already present} in \hohh$^+$
\book{150--151}:
\begin{itemize}
\item composing code is \emph{conjunction} of program clauses;
\item visibility of names is \emph{existential quantification} over program clauses (in effect, $\forall$
  in goals);
\item using a module to answer a query is an \emph{implicational goal} $E\imp G$.
\end{itemize}
The modules language is designed to support abstraction, information hiding, parametrization,
interfaces and separate compilation. It must not complicate the meaning of the core language, and it must
let higher-order programming coexist with modules (in Prolog, getting \texttt{call/1} to interact
correctly with modules). Finally, it should support a notion of module equivalence (representation
independence) \book{150--151}.

\subsection{Modules, signatures and accumulation}

A module is a named sequence of kind, type and clause declarations, and a signature with the same name is
its interface \book{152--153}. The book's module \lp{smlists} declares and defines \lp{id},
\lp{memb}, \lp{member} (a deterministic membership test that uses the cut), an auxiliary
\lp{revapp}, \lp{reverse} (via \lp{revapp}), \lp{append} and \lp{memb_and_rest} (select an element and
return the rest). Its signature is
\begin{lprolog}
sig smlists.
type id              list A -> list A -> o.
type memb, member    A -> list A -> o.
type reverse         list A -> list A -> o.
type append          list A -> list A -> list A -> o.
type memb_and_rest   A -> list A -> list A -> o.
end
\end{lprolog}
The signature restricts the external visibility of the module's names to exactly those it mentions, so
\lp{revapp} is hidden. Hiding applies to predicates and also to constants and data constructors, and
``such hidden symbols remain hidden even through the results of computations'' \book{153}. A query
\lp{[smlists] ?- revapp ...} is \emph{ill formed}: a static error, not a run-time failure \book{154}.
\lp{accumulate m1, ..., mn.} conceptually inserts the code of the listed modules, after \emph{renaming}
their private constants, so that the accumulated copy is entirely local to the accumulating module
\book{154}. \lp{accum_sig} inserts signatures, which is how names are re-exported \book{155--156}.

A module \emph{matches} a signature when its implicit signature (its own declarations, plus the
\emph{explicit} signatures of the modules it accumulates) is \emph{mergeable} with it: kinds are
identical, types are identical up to renaming of type variables, and operator declarations agree
\book{156--158}. The book's running example \book{158}:
\begin{lprolog}
module m1.
kind   item   type.
type   p,q    item -> o.
type   a      item.
type   r      list item -> o.
p X :- q X.
r (a :: nil).
end

sig m1.
kind   item   type.
type   p,q    item -> o.
type   a      item.
end
\end{lprolog}
\begin{lprolog}
module m2.
kind   item   type.
type   q, r   item -> o.
type   a      item.
q X :- r X.
r a.
end

sig m2.
kind   item   type.
type   q, r   item -> o.
type   a      item.
end
\end{lprolog}
\begin{lprolog}
module m3.
accumulate m1, m2.
type   s, t   item -> o.
type   b      item.
s X :- p X.
t b.
end

sig m3.
kind   item   type.
type   s, t   item -> o.
type   a      item.
end
\end{lprolog}
\lp{r} is hidden in \lp{m1} at type \lp{list item -> o}, while \lp{m2} exports an \lp{r : item -> o}. The
two never meet. If \lp{sig m1} also exported \lp{r}, the signatures of \lp{m1} and \lp{m2} would not be
mergeable and \lp{m3} would be ill formed \book{159--160}. Hiding is what lets two modules use the same
name at incompatible types. The name \lp{q}, on the other hand, is exported by both \lp{m1} and
\lp{m2} at the same type, so in \lp{m3} it is \emph{one} constant. It is the only link between
\lp{m1}'s clause \lp{p X :- q X} and \lp{m2}'s clause \lp{q X :- r X}.

\subsection{The logical reading: \texorpdfstring{$\exists$}{exists} over clauses, \texorpdfstring{$\imp$}{implication} for queries}\label{sec:hoprog-eformulas}

Two equivalences, valid in most logics, motivate the construction \book{160}:
\[
\exists y\,\forall x\,(D(x)\imp G(y)) \equiv \exists y\,((\exists x\,D(x))\imp G(y))
\qquad
\exists x\,(D_1(x)\wedge D_2) \equiv (\exists x\,D_1(x))\wedge D_2
\]
(when $x$ is not free in $G(y)$, respectively in $D_2$). The first says that a universal in a goal makes
$x$ \emph{local to the clause} $D(x)$. The second lets the scope of such an $\exists$ be narrowed to part of
a large conjunction.

\begin{definition}[E-formulas {\book{161}}]\label{def:hoprog-eformula}
The goals of \hohh$^+$ are extended to
\begin{align*}
G &::= \top \mid A \mid G\wedge G \mid G\vee G \mid \exists x\,G \mid E\imp G \mid \forall x\,G\\
D &::= A_r \mid G\imp D \mid D\wedge D \mid \forall x\,D\\
E &::= D \mid E\wedge E \mid \exists x\,E
\end{align*}
A rigid head may now also be a variable bound by an $\exists$ of an enclosing E-formula. E-formulas are
not user syntax. They appear only in the translation of modules \book{162}.
\end{definition}

A program contains only $D$-formulas, so ${\imp}R$ is replaced by a small phase that decomposes the
antecedent. It uses a sequent $\hoprogE{\Sig}{\Prog}{\Delta}{G}$, where $\Delta$ is a multiset of E-formulas
\book{161}:
\begin{mathpar}
\inferrule{\hoprogE{\Sig}{\Prog}{E}{G}}{\seq{\Sig}{\Prog}{E\imp G}}\;\rulename{$\imp$R$'$}
\and
\inferrule{\seq{\Sig}{\Prog}{G}}{\hoprogE{\Sig}{\Prog}{}{G}}\;\rulename{finish}
\and
\inferrule{\hoprogE{\Sig}{\Prog,D}{\Delta}{G}}{\hoprogE{\Sig}{\Prog}{D,\Delta}{G}}\;\rulename{reclassify}
\and
\inferrule{\hoprogE{\Sig}{\Prog}{E_1,E_2,\Delta}{G}}{\hoprogE{\Sig}{\Prog}{E_1\wedge E_2,\Delta}{G}}\;\rulename{$\wedge$L$'$}
\and
\inferrule{\hoprogE{\Sig,y}{\Prog}{E[y/x],\Delta}{G}}{\hoprogE{\Sig}{\Prog}{\exists x\,E,\Delta}{G}}\;\rulename{$\exists$L$'$}\ (y\notin\Sig)
\end{mathpar}
All these rules are \emph{invertible}, so this phase never backtracks. ``During this new phase, the local
constants are introduced dynamically into computation, and the clauses using these new constants are
integrated with the ambient logic program'' \book{161}. Note that $\exists L'$ is the same mechanism as
$\forall R$ for goals: a fresh constant is added to $\Sig$.

\begin{definition}[A module as a formula {\book{162}}]\label{def:hoprog-modformula}
Let $h_1,\ldots,h_k$ be the hidden constants of a module $M$, that is, those in its implicit signature but
not in its explicit one. If $M$ accumulates $M_1,\ldots,M_n$ and has its own clauses $C_1,\ldots,C_m$, then
\[
\hoprogT{M} = \exists h_1\cdots\exists h_k\,\big(\hoprogT{M_1}\wedge\cdots\wedge\hoprogT{M_n}\wedge C_1\wedge\cdots\wedge C_m\big).
\]
A query \lp{[M] ?- g X.} is the goal $\exists x\,(\hoprogT{M}\imp g\,x)$ \book{163}.
\end{definition}

For the running example \book{162--163}:
\begin{align*}
\hoprogT{m1} &= \exists r\,\big((\forall x(q\,x\imp p\,x))\wedge (r\,(a::nil))\big)\qquad
\hoprogT{m2} = \forall x(r\,x\imp q\,x)\wedge(r\,a)\\
\hoprogT{m3} &= \exists p\exists q\exists r\exists b\Big(\exists r\big((\forall x(q\,x\imp p\,x))\wedge(r\,(a::nil))\big)\\
&\qquad\qquad\wedge(\forall x(r\,x\imp q\,x)\wedge(r\,a))\wedge(\forall x(p\,x\imp s\,x)\wedge(t\,b))\Big)
\end{align*}
The inner $\exists r$ (\lp{m1}'s private \lp{r}) shadows the outer one (\lp{m2}'s exported \lp{r}, hidden
again by \lp{m3}), as $\alpha$-conversion dictates. \lp{a}, \lp{s} and \lp{t} stay free, because they are
\lp{m3}'s exports.

\begin{tutorial}[Using a hidden constant versus exposing it]
The book reports \book{163--164}:
\begin{lprolog}
[m3] ?- sigma x\ t x.      % succeeds (x instantiated to the hidden constant b)
[m3] ?- t X.               % FAILS
\end{lprolog}
\lp{[m3] ?- p a.} is ill formed (\lp{p} is not in \lp{sig m3}), and \lp{[m3] ?- s a.} succeeds. The
difference between the first two lies in where the quantifiers sit.
\begin{itemize}
\item \lp{t X} is $\exists x\,(\hoprogT{m3}\imp t\,x)$. The logic variable $X$ is created \emph{first}, at a
  signature that does not contain \lp{b}. Then ${\imp}R'$ and $\exists L'$ introduce a fresh $b'$.
  Backchaining on \lp{t b'} needs $X := b'$, which is not a term over $X$'s signature, so the goal fails.
\item \lp{sigma x\ t x} is $\hoprogT{m3}\imp\exists x\,(t\,x)$. $\exists L'$ adds $b'$ to the signature
  \emph{before} the existential is processed, so $b'$ is an admissible witness and the goal succeeds.
\end{itemize}
In the book's words, this is ``the difference between using a constant in the course of searching for a
solution \ldots\ and attempting to expose the identity of this constant externally'' \book{164}.
\end{tutorial}

\begin{tutorial}[{The rules of the E-phase at work on \texttt{[m3] ?- s a.}}]
The query has no logic variables, so it is the goal $\hoprogT{m3}\imp s\,a$. Read the derivation
bottom-up, starting from the empty program:
\begin{enumerate}
\item ${\imp}R'$ moves $\hoprogT{m3}$ into the E-multiset: $\hoprogE{\Sig}{\emptyset}{\hoprogT{m3}}{s\,a}$.
\item $\exists L'$ four times introduces fresh constants $p',q',r',b'$ for \lp{m3}'s hidden names and
  adds them to $\Sig$.
\item $\wedge L'$ splits the body into its conjuncts: \lp{m1}'s formula
  $\exists r\,(\ldots)$, the clauses $\forall x(r'\,x\imp q'\,x)$ and $r'\,a$ from \lp{m2}, and the
  clauses $\forall x(p'\,x\imp s\,x)$ and $t\,b'$ of \lp{m3}.
\item $\exists L'$ once more introduces a fresh $r''$ for \lp{m1}'s private \lp{r}, and $\wedge L'$
  yields $\forall x(q'\,x\imp p'\,x)$ and $r''\,(a::nil)$.
\item \rulename{reclassify} moves each of these $D$-formulas into the program, and \rulename{finish}
  returns to an ordinary sequent $\seq{\Sig,p',q',r',b',r''}{\Prog_{m3}}{s\,a}$.
\item Backchaining now proves $s\,a$ from $p'\,a$, $p'\,a$ from $q'\,a$, $q'\,a$ from $r'\,a$, and
  $r'\,a$ is a fact.
\end{enumerate}
Steps 1--5 make no choices, which is what invertibility means in practice: loading a module never
backtracks. The fact $r''\,(a::nil)$ is never used, and it could not be: no clause outside \lp{m1}'s
own formula can mention $r''$.
\end{tutorial}

\paragraph{Elaboration.} Rather than replay \cref{def:hoprog-modformula} for every query, an implementation
inlines the accumulated modules and \emph{renames} their private names. The elaborated \lp{m3} contains
the clauses \lp{p X :- q X.}, \lp{r' (a :: nil).}, \lp{q X :- r X.}, \lp{r a.}, \lp{s X :- p X.} and
\lp{t b.}, where \lp{r'} is a new name for \lp{m1}'s local \lp{r} \book{165}. The book's implementation
section adds that, at run time, the machinery that already handles alternating $\exists/\forall$ in
goals also handles local constants, provided they are annotated properly when the module is used to
prove queries \book{172}.

\subsection{Abstract datatypes}

Exporting a type constructor while hiding its data constructors gives an abstract datatype \book{166}:
\begin{lprolog}
sig stack.
kind store           type -> type.
type init            store A -> o.
type add, remove     A -> store A -> store A -> o.
end
\end{lprolog}
\begin{lprolog}
module stack.
kind store           type -> type.
type emp             store A.
type stk             A -> store A -> store A.
type init            store A -> o.
type add, remove     A -> store A -> store A -> o.
init emp.
add    X S (stk X S).
remove X (stk X S) S.
end
\end{lprolog}
\lp{[stack] ?- init A.} \emph{fails}, because it tries to exhibit the representation of the empty stack.
\lp{[stack] ?- sigma A\ sigma B\ sigma C\ init A, add 1 A B, remove X B C.} uses the representation
anonymously and succeeds with \lp{X = 1} \book{166}. A client such as \lp{graph_search} can hold store
values in its variables, but only \lp{init}, \lp{add} and \lp{remove} can examine them \book{167}.

\paragraph{Extension and parametrization.} An exported predicate can be \emph{extended} by later modules.
A propositional prover \lp{proplogic} exports the type \lp{form} and the predicate \lp{prove}, and a
module \lp{quantlogic} that accumulates it adds both constructors and clauses \book{168--169}:
\begin{lprolog}
sig quantlogic.
accum_sig proplogic.
kind term        type.
type all, some   (term -> form) -> form.
end

module quantlogic.
accumulate proplogic, smlists.
kind term        type.
type all, some   (term -> form) -> form.

prove L F :- memb_and_rest (some P) L L',
             pi c \ prove ((P c)::L) F.
prove L (all P) :- pi c\ prove L (P c).
prove L (some P) :- prove L (P T).
prove L F :- memb_and_rest (all P) L L',
             append L' [all P] L'', prove ((P T)::L'') F.
end
\end{lprolog}
This works because \lp{form} and \lp{prove} are in \lp{proplogic}'s \emph{explicit} signature, so
accumulation neither renames nor hides them. The recursive calls inside \lp{proplogic}'s clauses reach
the new quantifier clauses too. Both provers accumulate \lp{smlists}, which gives two private copies. The
book calls this ``at least wasteful'' \book{169} and offers an alternative: accumulate the
\emph{signature} \lp{smlists} into both modules, and supply one implementation when everything is
accumulated at the top \book{170}. Parametrization is then done by shared names that are closed at the
top, not by passing arguments explicitly.

\subsection{Higher-order programming and predicate visibility}\label{sec:hoprog-callp}

The test case for the second desideratum \book{171}:
\begin{lprolog}
sig comblibrary.
type  call  o -> o.
end

module comblibrary.
type  call  o -> o.
type  p  list int -> o.
p (1 :: nil).
call Q :- Q.
end

sig test.
type  test  list int -> o.
end

module test.
accumulate comblibrary.
type  test  list int -> o.
type  p  list int -> o.
p (2 :: nil).
test X :- call (p X).
end
\end{lprolog}
For \lp{[test] ?- test X.}, which \lp{p} does \lp{call (p X)} invoke: the one visible where \lp{call} is
defined (\lp{X = 1 :: nil}), or the one visible where the name is written (\lp{X = 2 :: nil})? ``The second
interpretation is the commonly used resolution'', and it ``is a natural consequence of the semantics
that we have presented'' \book{171}. In the elaborated module, \lp{comblibrary}'s \lp{p} is renamed
\lp{p'}, so the only answer is \lp{2 :: nil} \book{171--172}. \lp{p X} is a \emph{term} of type
\lp{o} whose head is fixed when the program is parsed. It is not a name looked up at run time.

\begin{kanren}[Modules: what Scheme scoping gives, and what it does not]
Scheme's lexical scope already does module elaboration with renaming: a local definition is private, and
each evaluation of a scope creates fresh closures. Goals are closures, so Figure~6.10 resolves to the
book's ``second interpretation'' automatically. A Scheme-level abstract type is equally easy:
\begin{scheme}
;; Figure 6.10: a "module" is a lexical scope returning its exported relations
(define call°                            ; module comblibrary, exports call
  (let ()
    (define-relation (p l) (=° l '(1)))  ; hidden
    (define (call g) g)
    call))
(define test°                            ; module test, exports test
  (let ()
    (define-relation (p l) (=° l '(2)))
    (define-relation (test x) (call° (p x)))
    test))
(show 1 (μkanren-run (q -1 #t) (test° q)))
;; Figure 6.6: the stack "abstract datatype", constructors made by gensym
(define-values (init° add° remove°)
  (let ((emp (gensym 'emp)) (stk (gensym 'stk)))
    (define-relation (init s) (=° s emp))
    (define-relation (add x s s*) (=° s* (list stk x s)))
    (define-relation (remove x s s*) (=° s (list stk x s*)))
    (values init add remove)))
(show 2 (μkanren-run (q -1 #t)
          (fresh° (a b c) (init° a) (add° 1 a b) (remove° q b c))))
(show 3 (μkanren-run (q -1 #f) (init° q)))
\end{scheme}
\begin{answers}
1 ((2))
2 (1)
3 ((λ () (quote emp223)))
\end{answers}
Line 1 is the book's answer, and line 2 uses the stack anonymously, as the \lp{sigma} query does. Line 3
shows the honest difference. Scheme hides the \emph{name} \scm{emp}, but the answer to \lp{init A} still
\emph{exhibits} the representation, where λProlog fails. Matching the book needs hidden constructors that
are eigen-constants, introduced when the module is used, and a scope check in unification that stops a
query variable created earlier from being bound to them. That is the same machinery as $\forall$ in
goals (\cref{sec:hoprog-hohhplus}). Two more things are not free. One is \emph{extensible} predicates
(\lp{quantlogic} adding clauses to \lp{proplogic}'s \lp{prove}), because a \scm{define-relation} is
closed. The other is parametrization by shared names closed at the top. Signature matching is mostly
about types and kinds, which an untyped miniKanren can only approximate (names and arities). The proposal
is in \cref{ch:design,ch:module}. One rule follows directly from Figure~6.10: never resolve a relation from
a symbol at run time, because that reintroduces exactly the Prolog \texttt{call/1} ambiguity.
\end{kanren}

% ======================================================================================================
\section{What a Scheme-hosted miniKanren gets for free}\label{sec:hoprog-free}

\Cref{tab:hoprog-free} collects the kanren boxes of this chapter.

\begin{table}[ht]
\centering\small
\caption{Book features versus a miniKanren whose relations are Scheme procedures.}\label{tab:hoprog-free}
\begin{tabular}{@{}>{\raggedright\arraybackslash}p{0.40\textwidth}>{\raggedright\arraybackslash}p{0.17\textwidth}>{\raggedright\arraybackslash}p{0.36\textwidth}@{}}
\toprule
Feature & With closures & What is needed beyond closures\\
\midrule
\texttt{mappred}, \texttt{foreach}, \texttt{forsome}, \texttt{sublist}, folds, \texttt{compose} & free & nothing (curried \texttt{define} for partial application)\\
$\lambda$-terms of predicate type as arguments & free & nothing (Scheme \texttt{λ})\\
closures \texttt{ref}/\texttt{sym}/\texttt{trans}, binary clauses (CPS) & free & nothing\\
hidden predicates with fixed clauses (\texttt{reverse}) & free & lexical scope\\
\texttt{if}, \texttt{not} & free & \texttt{if°} with \texttt{take° 1}\\
\texttt{mapfun}/\texttt{reducefun} forwards, \texttt{collect'} forwards & free & closures that build terms; \texttt{project°}\\
module hiding, predicate visibility (Fig.~6.10) & free & lexical scope\\
\midrule
\texttt{rel R, R john mary} & opaque answers & relations as terms, to print or compare them\\
flexible goals & accidental error & inspectable goals; a store of suspended goals\\
growing hidden predicates (\texttt{fib\_memo}), $\forall$/$\imp$ encodings of $\vee$, $\exists$ & no & goal-level $\forall$ (eigenvariables, scope check) and $\imp$ (dynamic clauses)\\
abstract datatypes that stay hidden in answers & partial & eigen-constants and a scope check\\
extensible predicates across modules & no & open clause sets, linking\\
\texttt{eq\_pred}, reasoning by substitution & no & predicates and clauses as $\lambda$-terms\\
\texttt{mapfun}/\texttt{reducefun} backwards, \texttt{extract\_a}, \texttt{palindrome}, fdl destructuring & no & $\lambda$-terms, $\beta\eta$-normalisation, higher-order unification\\
\bottomrule
\end{tabular}
\end{table}

The dividing line is simple. Whatever only \emph{calls} a predicate or a function is free, because a
Scheme procedure can be called. Whatever must \emph{look inside} one (unify it, print it, compare it,
synthesise it) needs $\lambda$-terms in the term language. A third group, $\forall$ and $\imp$ in goals,
needs neither $\lambda$-terms nor inspection but does need new components in the search state:
eigenvariables with a scope check, and a dynamic clause store. \Cref{ch:design} builds on this
classification.

% ======================================================================================================
\section{Summary}\label{sec:hoprog-summary}
\begin{itemize}
\item Atoms are \emph{rigid} (constant head) or \emph{flexible} (variable head). Clause heads must be rigid,
  which keeps programs consistent and modular; flexible atoms are allowed in goals.
\item Implications may not occur inside terms. Instantiation terms are drawn from Herbrand universes:
  $\mathcal{H}_1$ (no $\forall,\imp$) for \hohc{} and $\mathcal{H}_2$ (no $\imp$) for \hohh{}. Goal-directed
  search is then complete: classical and intuitionistic for \hohc{}, intuitionistic for \hohh{}.
\item \hohh$^+$ lets a variable bound by an essentially universal quantifier head assumed clauses, which
  gives hidden, closed predicate definitions.
\item Predicate variables give \lp{mappred}, \lp{sublist}, closures, \lp{compose}, \lp{foldl} and the CPS
  transformation. Logic variables inside a $\lambda$ are shared.
\item An unbound flexible goal is underconstrained. Suspending it is complete; in practice it is reported
  as an error. Constrain predicate variables first (\lp{rel R, R john mary}).
\item Hidden predicates support reasoning by substitution (the symmetry of \lp{reverse}). Logical
  constants are definable. \lp{if} and \lp{not} via cut are non-logical.
\item Function variables are evaluated by $\beta\eta$ and can be synthesised by unification
  (\lp{mapfun}, \lp{reducefun}, functional difference lists). Higher-order unification is not unitary, so
  prefer structural recursion and the right quantifier order.
\item Unlike functions in functional programs, predicate expressions can be compared
  intensionally, with \lp{eq_pred}.
\item Modules are logic: conjunction composes code, $\exists$ over clauses hides names, and a query
  against a module is $\exists x\,(\hoprogT{M}\imp G)$. Hidden constants may be used but not exposed, and
  predicate names resolve statically.
\item In a Scheme-hosted miniKanren, calling is free and looking inside is not; $\forall$ and $\imp$ in
  goals need new state (\cref{tab:hoprog-free}, \cref{ch:design}).
\end{itemize}
